
\medskip

Pour chacune des affirmations suivantes, dire si elle est vraie ou fausse en justifiant soigneusement la réponse. 

\medskip
\begin{enumerate}
\item 
\begin{minipage}{8cm}
La récolte de la lavande débute lorsque les trois quarts des fleurs  au moins sont fanées. Le producteur a cueilli un échantillon de lavande représenté par le dessin ci-contre.

\smallskip

\textbf{Affirmation 1 :} la récolte peut commencer.

\end{minipage}
\hfill
\begin{minipage}{6cm}
\begin{pspicture}(6,4)
%\psgrid
\rput(4.53,2){\includegraphics[width=2.9cm]{lavande}}
\uput[r](0,3.4){Fleur fraîche}
\uput[r](0,2.73){Fleur fânée}
\psline{->}(2.2,3.4)(3.6,3.4)
\psline{->}(2.2,2.73)(3.6,2.73)
\end{pspicture}
\end{minipage}

\item En informatique, on utilise comme unités de mesure les multiples de l’octet :

\[ 1 \text{ko} = 10^3 \text{octets},\quad 1 \text{Mo} = 10^6 \text{octets},\quad 1 \text{Go} = 10^9 \text{octets}. \]

\parbox{0.48\linewidth}{
Contenu du disque dur externe :
 
$\bullet~~$ \np{1000} photos de $900$ ko chacune: 

$\bullet~~$ 65 vidéos de 700 Mo chacune.}\hfill
\parbox{0.5\linewidth}{ 
Capacité de l’ordinateur : 250 Go 

\begin{pspicture}(4,2)
\psset{unit=1.5cm,hatchwidth=0.5pt,hatchsep=.1pt}
\scalebox{.99}[0.3]{\psellipticwedge[fillstyle=hlines,linewidth=2pt](1,3)(1,0.9){145}{90}}
\psline[linewidth=1.25pt](1,1.17)(1,0.9)(0.17,1.02)
\psline[linewidth=1.25pt](1.98,0.7)(1.98,0.9)
\psline[linewidth=1.25pt](0.01,0.7)(0.01,0.9)
\scalebox{.99}[0.3]{\psarc[linewidth=1.5pt](1,2.48){1}{180}{0}}%
\uput[r](2,1){\psframe[fillstyle=solid,fillcolor=black](0.2,0.2)\hspace{0.4cm} Espace utilisé : 200 Go}
\uput[r](2,0.5){\psframe(0.2,0.2)\hspace{0.4cm} Espace libre}
\scalebox{.99}[0.3]{\psarc[linewidth=2pt](1,3){0.9}{90}{145}}
\end{pspicture}
}

\smallskip

\textbf{Affirmation 2 :} le transfert de la totalité du contenu du disque dur externe vers l'ordinateur n'est pas possible.
\item On considère le programme de calcul ci-dessous :

\begin{center} 
\begin{tabularx}{0.5\linewidth}{|X|}\hline
Choisir un nombre ; \\
Ajouter 5 ;\\ 

Multiplier le résultat obtenu par 2 ;\\

Soustraire 9.\\ \hline
\end{tabularx}
\end{center} 

\textbf{Affirmation 3 :} ce programme donne pour résultat la somme de 1 et du double du nombre choisi. 
\end{enumerate} 

\vspace{0.5pt}

